\title{DiffAccel: Accelerating Diffusion Models Through Adaptive Feature Optimization and Dynamic Hardware Adaptation}

\begin{abstract}
Diffusion Models (DMs) have emerged as a revolutionary technology in AI-generated content. Despite their superior quality, DMs face significant deployment challenges due to two key limitations: 1) the substantial computational cost inherent in their iterative denoising process, and 2) the complexities of optimizing the execution of their multi-scale U-Net architecture, which exhibits diverse computational patterns and varying memory requirements across layers. Existing optimization approaches face limitations on both algorithmic and hardware fronts. At the algorithm level, current methods struggle with inefficient feature enhancement through uniform channel processing and inflexible feature reuse with static caching strategies. On the hardware side, current solutions lack the flexibility and efficiency to adapt to the varied computational and memory demands within the U-Net. To address these challenges, we propose DiffAccel, an algorithm-hardware co-optimized accelerator designed to accelerate the inference of various DMs. Firstly, two key algorithmic innovations are introduced: selective feature enhancement targeting only high-dimensional channels and an adaptive feature caching mechanism based on identified evolution patterns. Secondly, at the hardware level, a reconfigurable array is proposed to enable efficient execution across diverse U-Net layers by switching its dataflow between output-stationary and weight-stationary modes based on the specific computational patterns of each layer. Moreover, a data-volume-aware memory management unit is implemented, employing adaptive input allocation and hierarchical output storage to efficiently manage the disparate layer memory footprints of the U-Net, thereby improving on-chip memory efficiency and reducing waste. Thirdly, fine-grained pipeline fusion is employed for key U-Net operation sequences to minimize intermediate data movement and significantly reduce latency. Specifically, row-first processing is applied for linear-softmax operations, and channel-priority output is utilized for convolution-groupnorm sequences. At the algorithm level, DiffAccel achieves 2.93$\times$ speedup and reduces computational complexity (from 105 TFLOPs to 69 TFLOPs) compared to the baseline Stable Diffusion v2.0. DiffAccel is implemented with a 28nm technology. In terms of hardware efficiency, 
\textcolor{black}{based on post place-and-route implementation results, DiffAccel achieves 5.25$\times$-5.79$\times$ higher throughput and maintains comparable area efficiency} 
compared to state-of-the-art accelerators.
\end{abstract}

\begin{IEEEkeywords}
Diffusion Models, Feature Optimization, Hardware Adaptation.
\end{IEEEkeywords}

\section{Introduction}
% \IEEEPARstart{T}{his} file is intended to serve as a ``sample article file''
\IEEEPARstart{D}{iffusion} Models (DMs) have demonstrated exceptional creative capabilities in artificial intelligence-generated content (AIGC), showing remarkable performance in text-to-image generation \cite{text-img-1}\cite{text-img-2}, image super-resolution \cite{img-super-resolution-1}\cite{img-super-resolution-2}, image editing \cite{img-edit-1}\cite{img-edit-2}, and video generation \cite{video-gen-1} \cite{video-gen-2}. Their fundamental principle relies on a probabilistic, gradual denoising process that precisely transforms random noise into high-quality images. Compared to traditional Generative Adversarial Networks (GANs) \cite{gan-1}\cite{gan-2} and Variational Autoencoders (VAEs) \cite{vae-1}\cite{vae-2}, DMs deliver superior detail fidelity and generation quality.

However, the remarkable performance of DMs comes at a substantial computational cost. Unlike GANs, which require only a single forward pass for inference, DMs rely on an iterative denoising mechanism based on a Markov chain \cite{markov}. A standard diffusion process typically involves 50–100 denoising iterations. The computational cost is further exacerbated by the U-Net architecture \cite{unet} used in each iteration, which contains approximately 800 million parameters in Stable Diffusion (SD) v1.5 \cite{sd} (occupying~3.1GB in FP32 precision) \cite{sda}. These combined factors result in intensive computational demands. Consequently, generating a single $256 \times 256$ image takes 2.56 seconds on a high-performance NVIDIA A100 GPU \cite{isscc24}, while the same task requires 3.5 hours on resource-constrained embedded processors \cite{sda}.

To address these computational challenges, various algorithmic optimization approaches have been explored. Early efforts focused on model compression techniques, including knowledge distillation \cite{progressive-distill}\cite{knowledge-distillation} and post-training quantization methods such as Post-Training Quantization for Diffusion models (PTQD) \cite{ptqd} and Q-Diffusion \cite{q-diffusion}. However, these methods often degrade generation quality, limiting their effectiveness for large-scale models like SD. Recent research has shifted toward training-free optimizations addressing feature reuse and enhancement. Feature reuse methods like DeepCache \cite{deepcache}, Faster Diffusion \cite{faster-diffusion}, and Temporal Attention Decomposition (TGATE) \cite{TGATE} exploit temporal consistency but use fixed caching strategies that fail to adapt to dynamic diffusion processes. Feature enhancement approaches like Free Lunch in Diffusion U-Net (FreeU) \cite{freeu} improve quality by modulating skip connections and backbone features, yet introduce computational redundancy through uniform enhancement. Most existing studies focus on individual optimization strategies in isolation, missing potential synergistic effects.

\begin{figure}[!t]
\centering
\includegraphics[width=3.0in]{fig/pdf/moti.pdf}
\caption{The analysis of U-Net denoising in Stable Diffusion. (a) Data size distribution by layer in U-Net (under FP16); (b) Dynamic demand for memory resources at each layer.}
\label{fig:moti}
\end{figure}

Despite these algorithmic advancements in the inference efficiency of DMs, U-Net remains a critical performance bottleneck in DMs, accounting for over 98\% of the total execution time \cite{sda}. However, efficiently optimizing U-Net is particularly challenging due to its multi-scale architecture, which processes features through downsampling and upsampling pathways.
First, different resolution blocks exhibit inverse data access patterns. Quantitative analysis, as shown in Fig. \ref{fig:moti}(a), reveals significant disparities in data size distribution across U-Net layers. In high-resolution computational blocks (HRCBs), large spatial features are processed with fewer channels. The feature data volume in these blocks can reach up to 2560KB (layer 1), necessitating efficient feature data reuse to minimize memory access overhead. Conversely, in low-resolution computational blocks (LRCBs), smaller spatial features are handled with substantially more channels. Here, the weight data volume becomes the primary bottleneck, reaching approximately 3200KB in layers 4-8, which demands efficient weight parameter reuse.
Second, memory requirements vary dramatically across layers. For instance, operating at FP16 precision in SD v1.5 with a $512 \times 512$ input resolution, the required on-chip memory footprint per computational block varies drastically, ranging from 2.5MB for high-resolution blocks down to just 0.16MB for low-resolution blocks. As illustrated in Fig. \ref{fig:moti}(b), this substantial disparity poses a significant challenge for memory allocation. Conventional fixed allocation strategies are rendered ineffective: maximum-based allocation (2.5MB) leads to resource wastage, while minimum-based allocation (0.16MB) fails to meet the requirements of high-resolution layers.

Existing hardware acceleration strategies improve computational efficiency but fail to fully address the structural challenges within U-Net architectures, as shown in Fig. \ref{fig:moti}. Some works based on timestep similarity overlook the inherent architectural heterogeneity of U-Net \cite{isscc24} \cite{cambricon-d} \cite{ditto}. Efforts to reduce sampling steps decrease the overall workload but also fail to resolve the internal structural imbalances of U-Net \cite{edgediff}. In contrast, recent work \cite{cim-unet}\cite{unet_cim_isscc} identifies variations in computational density but lacks tailored strategies for different network components. Most critically, these solutions generally employ static resource allocation strategies, making them unable to adapt to the dynamically shifting computational and memory demands within the U-Net architecture.

To address these challenges, we propose DiffAccel, an algorithm-hardware co-optimized accelerator that systematically accelerates the inference of various DMs. Our major contributions are as follows:
\begin{itemize}
\item{} At the algorithm level, we propose a framework that intelligently synergizes selective feature enhancement with adaptive feature caching. Quantitative analysis has been conducted to demonstrate that selective enhancement of high-dimensional features achieves better quality while reducing computational overhead. A distinctive three-phase evolution pattern is discovered during the diffusion process. Based on this pattern, an adaptive caching strategy is developed. This adaptive strategy pinpoints diffusion steps where feature reuse can be applied effectively, achieving a more balanced trade-off between generation quality and computational efficiency.

\item{} At the hardware level, an efficient hardware accelerator architecture is proposed to mitigate the computational and memory challenge during the inference of U-Net. A Reconfigurable Array (RA) is implemented to satisfy layer-specific processing requirements. Specifically, the RA operates in two typical modes: the Output-Stationary (OS) mode for high-resolution layers and the Weight-Stationary (WS) mode for low-resolution layers. The selection of the mode is determined by the corresponding computational requirements of each U-Net layer. To meet diverse memory requirements, a data-volume-aware Memory Management Unit (MMU) is well designed. The MMU effectively manages the disparate U-Net layer memory footprints by employing data-volume-based adaptive allocation for input data and a hierarchical on-chip/off-chip storage system for output data, significantly enhancing on-chip memory efficiency and reducing resource waste.

\item{} Specific consecutive operation sequences in U-Net are carefully scheduled using processing schemes specifically tailored to their distinct data access patterns. Specifically, for Matrix Multiplication (MM) followed by Softmax/LayerNorm (LNorm) \cite{attention} operations, a row-priority processing scheme is employed. Conversely, convolution operations followed by GroupNorm (GNorm) \cite{groupnorm} are implemented with channel-priority output. By tailoring the processing scheme to match the associated data access patterns of each operation sequence, this approach establishes an efficient pipeline, minimizing memory access overhead through reduced intermediate data movement and the elimination of costly data rearrangement.

\item{} Extensive experiments are conducted to demonstrate the effectiveness of DiffAccel. At the algorithm level, our optimized design reduces computational requirements from 105 TFLOPs in baseline SD v2.0 to only 69 TFLOPs, achieving a 2.93$\times$ speedup while maintaining comparable generation quality (with only a 0.56 increase in Fréchet Inception Distance (FID) score and a 0.06 improvement in Contrastive Language-Image Pre-training (CLIP) score compared to the baseline SD v2.0). At the hardware level, DiffAccel delivers significant performance improvements. Compared with cutting-edge accelerators in open literature, \textcolor{black}{based on post place-and-route implementation results, DiffAccel achieves 5.25$\times$-5.79$\times$ higher throughput and maintains comparable area efficiency.}
\end{itemize}